% !TeX root = ../I_interlock_architecture.tex
% Doc I master diagram.  Three fault sources, three actors, two MPS paths.
% Read top to bottom: the hardware layer (upper two thirds) acts in under a
% microsecond and is independent of software; the EPICS alarm tree and the SNL
% state machine (lower third) only observe and then tidy up.
%
% Layout rules used here:
%   * every box in a column shares that column's centre x, so the vertical
%     connectors between them are exactly vertical;
%   * the MPS paths run in lanes that touch no box -- A across the gap at the
%     actor row, B straight down the centre; the SPEAR MPS and orbit permits
%     come down the lane at x = 134 into Slot 5;
%   * the HVPS Remote I/O path runs down the far right at x = 203 so that it
%     does not pass through the VXI Slot 5 box.
\begin{tikzpicture}[font=\scriptsize, x=1mm, y=1mm]
\def\xA{30}\def\xB{100}\def\xC{168}
\def\yE{-150}

% ===== fault sources =======================================================
\node[blk,text width=54mm,anchor=north] (s1) at (\xA,-2)
  {Arc\\ {\tiny fiber}};
\node[blk,text width=58mm,anchor=north] (s2) at (\xB,-2)
  {Collector, vacuum, cooling,\\ reflected power {\tiny (protection)}};
\node[blk,text width=54mm,anchor=north] (s3) at (\xC,-2)
  {HVPS internal faults\\ {\tiny oil, arc, voltage, crowbar}};
\draw[grp] ($(s1.north west)+(-5,7)$) rectangle ($(s3.south east)+(5,-3)$);
\node[lbl,anchor=north west] at ($(s1.north west)+(-3,6)$) {FAULT SOURCES};

% ===== the three actors ====================================================
\def\yA{-28}
\node[rf,text width=54mm,anchor=north] (fic) at (\xA,\yA)
  {\textbf{FAST INTERLOCK}\\\textbf{CHASSIS 340-308}\\
   B132, $<\qty{1}{\micro\second}$\\[2pt]
   {\tiny photodetectors, peak detectors\\ no CPU, no firmware}};
\node[plc,text width=58mm,anchor=north] (mps) at (\xB,\yA)
  {\textbf{RF MPS PLC}\\ AB PLC-5 with 1771 I/O, B132\\
   $\approx\qty{10}{\milli\second}$ scan cycle};
\node[plc,text width=54mm,anchor=north] (slc) at (\xC,\yA)
  {\textbf{SLC-500 HVPS PLC}\\ B118, $\approx\qtyrange{10}{20}{\milli\second}$ scan\\[2pt]
   {\tiny detects arc, crowbar, overvoltage,\\ oil, transformer}};
\foreach \s/\d in {s1/fic, s2/mps, s3/slc}{\draw[sig] (\s.south) -- (\d.north);}

% ===== hardware outputs at each end ========================================
\node[pwr,text width=50mm,anchor=north] (fout) at (\xA,-60)
  {ARC SIGNALS {\tiny($\rightarrow$ AIM Slot 12)}\\
   AIM commands {\tiny($\rightarrow$ RF Local Display Panel)}\\
   {\tiny panel logic: SCR ENABLE, CROWBAR $\rightarrow$ B118}};
\draw[sig] (fic.south) -- (fout.north);

\node[pwr,text width=54mm,anchor=north] (sout) at (\xC,-60)
  {SCR supervisory enable\\ Slot 5 OX8 output 0 (24 V)\\
   {\tiny COMMANDS bus of the trigger boards, B118}};
\draw[sig] (slc.south) -- (sout.north);

% ===== Paths A and B =======================================================
% The actor boxes differ in height, so their .east/.west anchors sit at
% different y.  Route from one anchor and project onto the other's x.
\draw[sig] (mps.west) -- node[above,tag]{\textbf{Path A}} (mps.west -| fic.east);
\draw[draw=cNeut,thick] (mps.south) -- (\xB,-50);
\node[tag,anchor=west] at (\xB+3,-53) {\textbf{Path B}};

% ===== AIM and Slot 5 ======================================================
\def\yM{-84}
\node[ctrl,text width=76mm,anchor=north] (aim) at (\xA,\yM)
  {\textbf{AIM MODULE} --- VXI Slot 12 \texttt{devP2RfAim.c}\\[2pt]
   {\tiny arc comparators and latches \mbox{\texttt{ARCCURSTT}}, \mbox{\texttt{ARCLTDSTT}},
    thresholds from \mbox{\texttt{TRIPLVL}}}\\[2pt]
   asserts RF\_FAULT on the VXI P2 backplane};
\draw[sig] (fout.south) -- node[right,tag,xshift=1mm]{37-pin cable, J2 to J4} (aim.north);

\node[ctrl,text width=60mm,anchor=north] (slot5) at (\xC,\yM)
  {\textbf{VXI SLOT 5} --- MPS SHUTOFF MODULE\\[2pt]
   {\tiny inputs: SPEAR MPS permit $\cdot$ orbit-interlock permit (24 V)\\
    J3 $\rightarrow$ Fast IC spare fault $\rightarrow$ AIM External MPS}\\[2pt]
   asserts RF\_FAULT on P2};

% The machine permits come down the clear lane at x = 134, between the MPS and
% HVPS columns; the last turn is projected from the Slot 5 anchor.
\node[tag,anchor=north,align=center] (perm) at (134,-46) {SPEAR MPS,\\ orbit permits};
\draw[sig] (perm.south) -- (perm.south |- slot5.west) -- (slot5.west);

% ===== the RFP module, cut by either asserter ==============================
\node[rf,text width=54mm,anchor=north] (rfp) at (\xA,-118)
  {\textbf{RFP MODULE} --- Slot 4\\ RF drive DAC $\rightarrow$ zero\\
   {\tiny hardware, $<\qty{1}{\micro\second}$}};
\draw[sig] (aim.south) -- (rfp.north);
% dropped from the left of Slot 5 so it never shares the far-right RIO lane
\draw[sig] ($(slot5.south west)+(16,0)$) |- node[pos=0.62,above,tag]
  {RF\_FAULT line (open-collector)} (rfp.east);

% ===== EPICS alarm tree ====================================================
\node[blk,text width=212mm,anchor=north] (epics) at (98,\yE)
  {\textbf{EPICS ALARM AGGREGATION TREE}\quad
   \texttt{rf\_sumy\_stn.db} / \texttt{rf\_sumy\_stn\_spr.db}\quad
   {\tiny($\approx\qty{100}{\milli\second}$ CA propagation)}\\[3pt]
   {\footnotesize\texttt{STN:VXI:LTCH} $\rightarrow$ \texttt{STNPARK:SUMY:STAT}
    \qquad \texttt{STN:ABSUMY:LTCH}, \texttt{STN:SUMY:PLC},
    \texttt{STN:LOCAL:ON}
    \qquad \texttt{HVPSOFF:SUMY:STAT} $\rightarrow$ \texttt{HVPSSTN:SUMY:LTCH},
    \texttt{HVPSCONTACT:SUMY:STAT}}\\[3pt]
   \textbf{\texttt{\{STN\}:STNOFF:SUMY:STAT.SEVR} --- the single SNL trip wire}};

% AIM ISR: the software notification path.  Dropped straight down from a point
% on the AIM box's own bottom edge -- targeting a fixed x instead tilts it.
\coordinate (aimdrop) at ($(aim.south west)+(8,0)$);
\draw[sig] (aimdrop) -- (aimdrop |- 0,\yE);
\node[tag,anchor=west,align=left] at ($(aimdrop |- 0,\yE)+(4,15)$)
  {AIM ISR ($\approx\qty{10}{\micro\second}$)\\ $\rightarrow$ \texttt{STN:VXI:LTCH} $=$ MAJOR};
\draw[sig] (\xB,-50) -- (\xB,\yE);
% the HVPS Remote I/O path keeps to the far right of the Slot 5 box
\draw[sig] (sout.east) -| (203,\yE);
\node[tag,anchor=east] at (201,-114) {RIO $\approx\qty{1}{\hertz}$};

% ===== SNL state machine ===================================================
\node[ctrl,text width=150mm,anchor=north] (snl) at (98,\yE-30)
  {\textbf{SNL STATE MACHINE} \texttt{rf\_states.st} --- VxWorks IOC,
   $\approx\qty{1}{\second}$ response\\[3pt]
   {\footnotesize \texttt{s\_go\_off}: force beam abort $\rightarrow$ HVPS setpoint to zero
    $\rightarrow$ \qty{5}{\second} ramp $\rightarrow$ SCR gate off $\rightarrow$ RF drive off
    $\rightarrow$ fault file capture (11 channels)}};
\draw[sig] (epics.south) -- node[right,tag,xshift=1mm]
  {\texttt{fault\_stnoff} $\neq$ NO\_ALARM} (snl.north);
\end{tikzpicture}
